\section{Task Implementation Details}
\label{app:tasks}

Each CogArena task is implemented as a standalone web application following a consistent file structure. The entry point is an \texttt{index.html} file that loads jsPsych v8 and the experiment script. The experiment logic resides in \texttt{experiment.js}, which defines the full trial sequence including instruction screens, practice blocks with feedback, experimental blocks, and automatic data submission. A \texttt{task\_config.json} file specifies all configurable parameters including the number of trials, number of blocks, response keys or slider range, timing parameters (fixation duration, stimulus duration, response deadline, feedback duration), and condition proportions. The scoring directory contains \texttt{level2\_metrics.json} specifying accuracy metrics with human baseline statistics, and \texttt{level3\_signatures.json} specifying behavioral signatures with their statistical tests and expected directions.

Stimulus randomization is seeded from the session identifier, ensuring that each evaluation session produces a unique but reproducible stimulus sequence. For the perception and attention tasks, trial orders and condition assignments are shuffled within blocks. For the Two-Armed Bandit task, reward means and information conditions are generated per game. For the Iowa Gambling Task, loss schedules follow the classic Bechara payoff structure but loss timing is randomized. For the Reversal Learning task, the reversal schedule is generated dynamically with reversals occurring every 15 to 25 trials.

\section{Full Behavioral Signature Specifications}
\label{app:signatures}

This appendix provides the complete specification for all 32 behavioral signatures. Each entry describes the cognitive phenomenon being tested, the statistical test employed, the data fields and filters used, and the expected direction derived from the established human literature.

The Stroop task evaluates four signatures. Stroop interference in reaction time tests whether incongruent trials produce slower responses than congruent trials via independent-samples $t$-test (weight 1.0). Stroop interference in accuracy tests whether congruent trials produce higher accuracy than incongruent trials via two-proportion $z$-test (weight 0.5). Post-error slowing tests whether an error on the previous trial predicts slower reaction time on the current trial via sequential linear regression with a lagged inverted accuracy predictor (weight 1.0). The congruency sequence effect tests whether the Stroop interference effect is reduced following incongruent trials via a permutation-based interaction test on condition and previous-trial condition (weight 0.75).

The Go/No-Go task evaluates three signatures. Above-chance discrimination tests whether the hit rate on go trials exceeds the false alarm rate on no-go trials (weight 1.0). The commission-over-omission pattern tests whether false alarms (commission errors on no-go trials) exceed misses (omission errors on go trials), consistent with the prepotent go response (weight 0.75). Post-error slowing tests the same lagged relationship as in the Stroop task (weight 0.5).

The Flanker task evaluates four signatures following the same structure as the Stroop task: reaction time interference between incongruent and congruent trials (weight 1.0), accuracy interference (weight 0.75), post-error slowing (weight 0.5), and the congruency sequence effect (weight 0.5).

The Two-Armed Bandit task evaluates three signatures. Horizon-dependent exploration tests whether agents explore more in longer-horizon games via paired $t$-test comparing exploration rates between horizon-6 and horizon-1 conditions (weight 1.0). Win-stay behavior tests whether agents are more likely to repeat a choice after receiving a reward, exceeding chance (weight 0.75). Directed exploration tests whether agents preferentially choose the less-sampled arm when information is unequal in the long-horizon condition (weight 1.0).

The Risky Choice task evaluates three signatures. Risk aversion in gains tests whether agents choose the safe option more often than chance in the gain domain (weight 1.0). Loss aversion framing tests whether risk-seeking is higher in the loss domain than in the gain domain (weight 1.0). Expected value sensitivity tests whether the probability of choosing the risky option correlates positively with the expected value difference between risky and safe options (weight 0.75).

The Iowa Gambling Task evaluates three signatures. The learning effect tests whether the proportion of advantageous deck choices is higher in the final block than in the first block (weight 1.0). Above-chance advantageous choices tests whether the overall rate of advantageous deck selection exceeds 50\% (weight 0.75). Deck A avoidance tests whether agents learn to avoid the most punishing deck (weight 0.5).

The Trust Game evaluates three signatures. Non-zero trust tests whether the proportion of rounds with positive investment exceeds chance (weight 1.0). Reciprocity sensitivity tests whether previous return rates positively correlate with subsequent investment amounts (weight 1.0). Trustee type adaptation tests whether investment is higher with cooperative than with defecting trustees (weight 0.75).

The Dictator Game evaluates three signatures. Non-zero giving tests whether the proportion of rounds with positive giving exceeds chance (weight 1.0). Giving consistency tests whether giving amounts are stable across rounds via correlation with trial index (weight 0.5). Moderate giving tests whether the non-zero giving rate exceeds a more stringent threshold (weight 0.75).

The N-Back task evaluates three signatures. Above-chance discrimination tests whether the hit rate exceeds the false alarm rate (weight 1.0). Lure susceptibility tests whether false alarm rates are higher on lure trials (items matching the $n$-1 or $n$+1 position) than on non-lure non-target trials (weight 0.75). Post-error slowing follows the same sequential regression structure as in the attention tasks (weight 0.5).

The Reversal Learning task evaluates three signatures. Above-chance pre-reversal accuracy tests whether accuracy exceeds 50\% during the stable pre-reversal phase (weight 1.0). The perseveration effect tests whether pre-reversal accuracy exceeds post-reversal accuracy, reflecting the tendency to persist with the previously rewarded stimulus (weight 0.75). Post-reversal learning tests whether accuracy improves with increasing trials since the most recent reversal, indicating recovery from the contingency change (weight 0.75).

\section{Level 2 Accuracy Metrics}
\label{app:metrics}

CogArena computes 40 accuracy metrics across ten tasks. Each metric is normalized against human baselines via $z$-score transformation followed by cumulative normal distribution mapping to produce a score between 0 and 1. The Stroop task contributes four metrics: overall accuracy (human $M = 0.95$, $\mathit{SD} = 0.04$), congruent accuracy ($M = 0.97$, $\mathit{SD} = 0.03$), incongruent accuracy ($M = 0.92$, $\mathit{SD} = 0.06$), and mean correct reaction time ($M = 680$ ms, $\mathit{SD} = 120$ ms, lower is better). The Go/No-Go task contributes five metrics: overall accuracy, go accuracy, no-go accuracy, signal detection $d'$ ($M = 3.0$, $\mathit{SD} = 0.8$), and mean go reaction time. The Flanker task contributes four metrics following the same structure as the Stroop task. The remaining seven tasks each contribute three to five metrics tailored to their specific paradigms, including proportion of advantageous choices (Iowa Gambling), mean giving proportion (Dictator Game), investment proportion (Trust Game), exploration rates by horizon (Bandit), and pre- and post-reversal accuracy (Reversal Learning).

\section{API Reference}
\label{app:api}

The CogArena REST API provides nine endpoints. \texttt{GET /api/info} returns server information including the list of available tasks and endpoint descriptions. \texttt{GET /api/health} returns a health check response. \texttt{GET /api/tasks} returns full task metadata including configuration parameters, trial counts, response types, and the number of behavioral signatures. \texttt{POST /api/sessions} creates a new evaluation session given an agent name, scaffold, and model identifier, and returns the session identifier along with task URLs. \texttt{GET /api/sessions/\{session\_id\}} returns the current session status and task completion progress. \texttt{POST /api/data/\{session\_id\}/\{task\_id\}} receives trial data as a JSON array of trial objects. \texttt{POST /api/evaluate/\{session\_id\}} triggers the scoring pipeline for all submitted tasks in the session. \texttt{GET /api/results/\{session\_id\}} returns the full scorecard including per-task Level~1, Level~2, and Level~3 scores, individual metric values, signature test results, and the composite CogArena Score. \texttt{GET /api/leaderboard} returns ranked agent results across all scored sessions.
